\Article{11}
Zur weiteren Erläuterung dieser Methode führen wir eines ihrer zahlreichen Anwendungsbeispiele an. Die am $1832$ Sept. $11$ angestellten Versuche ergaben folgende Zahlen.
\SourceLeaf{102}
\expandafter\def\csname HistoricalNoteText10\endcsname{So betrug die Schwingungsamplitude der ersten Nadel im ersten Versuch anfangs $0^0\ 37'26''$ und am Ende $0^0\ 28'34''$. Im fünften Versuch betrug sie anfangs $1^0\ 10'21''$, nach $177$ Schwingungen $0^0\ 45'35''$ und nach $677$ Schwingungen $0^0\ 6'44''$.}
\begin{center}
\begin{tabular}{c||c|c||c}
&\multicolumn{3}{c}{Gleichzeitige Schwingungen}\\
Vers.&\multicolumn{2}{c||}{erste Nadel}&zweite Nadel\\\cline{1-4}
&Belastung&eine Schwingung&eine Schwingung\\
I&$r=180^{\mathrm{mm}}$&$24''63956$&$17''32191$\\
II&$r=130$&$20{,}77576$&$17{,}32051$\\
III&$r=80$&$17{,}66798$&$17{,}31653$\\
IV&$r=30$&$15{,}80310$&$17{,}30529$\\
V&ohne Belastung&$15{,}22990$&$17{,}31107$
\end{tabular}
\end{center}

Die Zeiten wurden mit einem Chronometer gemessen, das im Laufe eines Tages mittlerer Zeit um $14''24$ nachging. Jedes der beiden Gewichte $p$ betrug $103{,}2572$ Gramm. Die Abstände $r$ wurden in Millimetern mit mikroskopischer Genauigkeit bestimmt. Die Dauer einer Schwingung wurde aus mindestens $100$ Schwingungen (im fünften Versuch für die erste Nadel sogar aus $677$) ermittelt und bereits auf die Dauer einer unendlich kleinen Schwingung zurückgeführt. Diese Umrechnungen sind übrigens wegen der sehr kleinen Schwingungsamplitude\HistoricalNote{10} unmerklich, die wir bei unseren Apparaten unter Wahrung höchster Genauigkeit verwenden können. Wir führen diese Schwingungszeiten zunächst auf den Mittelwert von $TM$ zurück, der während des fünften Versuchs vorlag, und bedienen uns dabei der Vorschriften des vorangehenden Artikels unter II. Dann führen wir sie durch Multiplikation mit $\BareRoot{\frac{n+1}{n}}$ auf die Werte zurück, die sich ohne Torsion ergeben hätten. Dabei ist $n$ in den ersten vier Versuchen $=424{,}8$, im fünften $=597{,}4$ (vgl. Artikel $9$). Schließlich führen wir sie durch Multiplikation mit $\frac{86400}{86385{,}76}$ auf die mittlere Sonnenzeit zurück. So erhalten wir
\begin{center}
\begin{tabular}{crl}
I.&$24''65717=t'$&für $r'=180^{\mathrm{mm}}$\\
II.&$20{,}79228=t''$&für $r''=130^{\mathrm{mm}}$\\
III.&$17{,}68610=t'''$&für $r'''=80^{\mathrm{mm}}$\\
IV.&$15{,}82958=t''''$&für $r''''=30^{\mathrm{mm}}$\\
V.&$15.24515=t$&für die unbelastete Nadel.
\end{tabular}
\end{center}

Nehmen wir die Sekunde, den Millimeter und das Milligramm als Einheiten für Zeit, Abstand und Masse, sodass $p=103257.2$ ist, so leiten wir aus der Verbindung des ersten mit dem vierten Versuch Folgendes ab.
\SourceLeaf{103}
\[
TM=179641070,\qquad K+C=4374976000
\]
und dann aus dem fünften Versuch
\[
K=4230282000,\quad \text{und}\quad C=144694000
\]
Will man jedoch alle Versuche in die Rechnung einbeziehen, so wendet man die Methode der kleinsten Quadrate am bequemsten folgendermaßen an. Wir gehen von den Näherungswerten der Unbekannten $x,y$ aus, die sich aus der Verbindung des ersten und vierten Versuchs ergeben. Wir bezeichnen die noch hinzuzufügenden Korrekturen mit $\xi,\eta$ und setzen
\[
\begin{aligned}
x&=88{,}13646+\xi\\
y&=21184{,}85+\eta
\end{aligned}
\]
So ergeben sich die berechneten Werte der Zeiten $t'.\ t'',t''',t''''$ mit den naheliegenden Methoden
\[
\begin{aligned}
t'&=24{,}65717-0{,}13988\xi+0{,}00023008\eta\\
t''&=20{,}78731-0{,}11793\xi+0{,}00027291\eta\\
t'''&=17{,}69121-0{,}10036\xi+0{,}00032067\eta\\
t''''&=15{,}82958-0{,}08980\xi+0{,}00035838\eta
\end{aligned}
\]
Ihr Vergleich mit den beobachteten Werten nach der Methode der kleinsten Quadrate liefert
\[
\begin{aligned}
\xi&=-0{,}03230,&\quad\eta&=-12{,}38\\
x&=88{,}10416,&y&=21172{,}47
\end{aligned}
\]

Hieraus ergibt sich schließlich
\[
TM=179575250,\qquad K+C=4372419000
\]
und dann mithilfe des ersten Versuchs
\[
K=4228732400,\qquad C=143686600
\]

Es folgt der Vergleich der aus den korrigierten Werten der Größen $x,y$ berechneten Zeiten mit den beobachteten Zeiten.
\begin{center}
\begin{tabular}{c|r|r|r}
Vers.&Berechnete Zeit&Beobachtete Zeit&Differenz\\\hline
I&$24''65884$&$24''65717$&$+0''00167$\\
II&$20{,}78774$&$20{,}79228$&$-0{,}00454$\\
III&$17{,}69046$&$17{,}68610$&$+0{,}00436$\\
IV&$15{,}82805$&$15{,}82958$&$-0{,}00153$
\end{tabular}
\end{center}

\SourceLeaf{104}
Die Länge des einfachen Pendels in Göttingen setzen wir $=994^{\mathrm{mm}}126$. Damit ist die Schwerkraft, gemessen in der den vorangehenden Rechnungen zugrunde liegenden Einheit der beschleunigenden Kräfte, $=9811{,}63$. Ziehen wir es vor, die Schwerkraft selbst als Einheit zu nehmen, so ergibt sich $TM=18302{,}29$. Diese Zahl gibt an, wie viele Milligramm unter der Wirkung der Schwerkraft auf einen ein Millimeter langen Hebel drücken müssen, um dieselbe Wirkung zu erzielen wie die Kraft, mit der der Erdmagnetismus jene Nadel um die vertikale Achse zu drehen sucht.

\Article{12}
Nachdem wir das Produkt aus der horizontalen erdmagnetischen Kraft $T$ und dem magnetischen Moment $M$ einer gegebenen Nadel bestimmt haben, gehen wir nun zum zweiten Teil der Untersuchung über, zur Bestimmung des Quotienten $\frac MT$. Dazu vergleichen wir die Wirkung dieser Nadel auf eine andere Nadel mit der Wirkung des Erdmagnetismus auf dieselbe Nadel. Wie schon in der Einleitung dargelegt, lässt sich dieser Vergleich sowohl im Zustand der Bewegung als auch im Gleichgewicht anstellen. Wir haben beide Methoden häufig erprobt. Da die letztere der ersteren aus mehreren Gründen bei Weitem vorzuziehen ist, beschränken wir unsere Untersuchung hier auf sie. Dies liegt umso näher, als sich die erstere ohne Schwierigkeit auf ganz dieselbe Weise behandeln lässt.

\Article{13}
Mit dem Prinzip der virtuellen Verschiebungen lassen sich die Gleichgewichtsbedingungen eines beweglichen Körpers, auf den beliebige Kräfte wirken, sehr leicht in einer einzigen Formel zusammenfassen. Für jede Kraft nimmt man die unendlich kleine Verschiebung ihres Angriffspunkts, projiziert sie in die Richtung dieser Kraft und multipliziert sie mit der Kraft. Die Summe dieser Produkte darf bei keiner virtuellen Verschiebung, also bei keiner Verschiebung, die mit den allgemeinen Bedingungen für die Bewegung des Körpers vereinbar ist, einen positiven Wert annehmen. Soweit daher virtuelle Verschiebungen in entgegengesetzten Richtungen überall möglich sind, muss diese Summe, die wir mit $\mathrm{d}\Omega$ bezeichnen, für jede virtuelle Verschiebung $=0$ sein.

Der bewegliche Körper, den wir hier betrachten, ist eine Magnetnadel, deren Punkt $G$ mit einem oben fixierten, tordierbaren Faden verbunden ist. Der Faden verhindert lediglich, dass der Abstand des Punkts $G$ vom festen Fadenende größer als die Fadenlänge wird. Somit hängt auch hier, wie bei einem völlig freien Körper, die Lage des Körpers im Raum von sechs Variablen und sein Gleichgewicht daher von sechs Bedingungen ab. Da die Lösung des Problems hier lediglich zur Bestimmung des Quotienten $\frac MT$ dienen \SourceLeaf{105}soll, genügt es, die virtuelle Verschiebung zu betrachten, die in der Drehung um die vertikale Achse durch $G$ besteht. Offenbar dürfen wir diese Achse als fest und allein den Winkel zwischen der vertikalen Ebene, in der die magnetische Achse der Nadel liegt, und der Ebene des magnetischen Meridians als veränderlich betrachten. Wir zählen diesen Winkel vom nördlichen Teil des Meridians nach Osten und bezeichnen ihn mit $u$.

